\section{Zadanie 10}
\autor{Mikołaj Kołek}

Startujemy doświadczenie z liczbą \(k\) (\(0 \leq k \leq n\)). Barierami są liczby \(0\) i \(n\). Dopóki aktualna liczba nie jest jedną z barier do aktualnej liczby dodajemy \(1\) lub \(-1\) z równym prawdopodobieństwem. Jakie jest prawdopodobieństwo, że spacer uderzy w którąś z barier? Jakie jest prawdopodobieństwo, że spacer uderzy w \(n\)?

\subsection*{Rozwiązanie}
Pierwsze rozwiążmy prostszą część, czyli prawdopodobieństwo uderzenia w dowolną z barier. Możemy zauważyć, że jeśli w naszym ciągu jedynek i minus jedynek pojawi się w jakiejkolwiek chwili \(n - 1\) jedynek pod rząd, to wiemy, że ten ciąg na pewno uderza w którąś z barier. Zadanie to sprowadza się więc do zadania 9, z którego wiemy, że prawdopodobieństwo wystąpienia takiego ciągu jedynek w naszym ciągu wynosi 1.

Teraz rozważmy prawdopodobieństwo, że spacer uderzy w \(n\). Niech:

\(X_i\) - zdarzenie polegające na tym, że spacer uderzył w \(n\) dla \(k = i\).

Można zauważyć, że \(P(X_0) = 0\), bo dla \(k = 0\) spacer natychmiastowo uderza w barierę \(0\), a \(P(X_n) = 1\), bo dla \(k = n\) spacer od razu uderza w \(n\). Dodatkowo możemy zauważyć, że:

$$ \forall_{i \in [n-1]} P(X_i) = \frac{P(X_{i-1})}{2} + \frac{P(X_{i+1})}{2} $$

Używając tego wzoru, możemy teraz policzyć stosunek \(P(X_1)\) do \(P(X_2)\), później \(P(x_2)\) do \(P(x_3)\), itd. Udowodnię indukcyjnie, że \(P(X_i) = \frac{i}{i+1} P(X_{i+1})\). Baza indukcji:

$$ P(X_1) = \frac{P(X_{0})}{2} + \frac{P(X_{2})}{2} = \frac{P(X_{2})}{2} $$

Krok indukcyjny:
$$ P(X_i) = \frac{i}{i+1} P(X_{i+1}) $$
$$ P(X_{i+1}) = \frac{P(X_{i})}{2} + \frac{P(X_{i+2})}{2} = \frac{iP(X_{i+1})}{2i +2} + \frac{P(X_{i+2})}{2}$$
$$ \frac{P(X_{i+1})(i+2)}{2i+2} = \frac{P(X_{i+2})}{2} $$
$$ P(X_{i+1}) = \frac{P(X_{i+2})}{2} \frac{2i+2}{i+2} = \frac{i+1}{i+2} P(X_{i+2}) $$

W takim razie, zachodzi:
$$ P(X_i) = \frac{i}{i+1} P(X_{i+1}) = \frac{i}{i+1}\frac{i+1}{i+2} P(X_{i+2}) = \frac{i}{i+1} \frac{i+1}{i+2} \ldots \frac{n-1}{n} P(X_n) = \frac{i}{n} P(X_n) = \frac{i}{n}  $$

A więc dla ustalonego \(k\), prawdopodobieństwo, że spacer uderzy w \(n\) to \(\frac{k}{n}\).

\subsection*{Odpowiedź}
\[
\boxed{\text{Prawdopodobieństwo, że spacer uderzy w którąś z barier = 1}}
\]
\[
\boxed{\text{Prawdopodobieństwo, że spacer uderzy w \(n\) = \(\frac{k}{n}\)}}
\]